\documentclass[aps,pra,reprint,superscriptaddress,nofootinbib,longbibliography]{revtex4-2}

\usepackage{amsmath}
\usepackage{amssymb}
\usepackage{amsthm}
\usepackage{bm}
\usepackage{graphicx}
\usepackage{hyperref}
\usepackage{orcidlink}

\newtheorem{lemma}{Lemma}

\newcommand{\bra}[1]{\langle #1 |}
\newcommand{\ket}[1]{| #1 \rangle}
\newcommand{\braket}[2]{\langle #1 | #2 \rangle}
\newcommand{\Tr}{\mathrm{Tr}}
\newcommand{\eps}{\varepsilon}
\newcommand{\HH}{\mathcal{H}}
\newcommand{\EE}{\mathcal{E}}
\newcommand{\SSys}{\mathcal{S}}
\newcommand{\E}{\mathbb{E}}
\newcommand{\PP}{\mathbb{P}}
\newcommand{\AAlg}{\mathcal{A}}
\newcommand{\MM}{\mathcal{M}}

\begin{document}

\title{The Geometric Part of Decoherence: \\ Quasi-Orthogonality in High-Dimensional Hilbert Spaces}

\author{Karl Svozil\,\orcidlink{0000-0001-6554-2802}}
\affiliation{Institute for Theoretical Physics, TU Wien,
Wiedner Hauptstrasse 8-10/136, A-1040 Vienna, Austria}
\email{karl.svozil@tuwien.ac.at}

\date{\today}

\begin{abstract}
We isolate a geometric mechanism that complements the dynamical suppression of macroscopic interference: In a high-dimensional Hilbert space, almost all state vectors are nearly orthogonal, accommodating an exponentially large reservoir of mutually quasi-orthogonal environmental records. This geometry explains why macroscopic alternatives fail to exhibit visible interference once such records are populated. The argument is conditional and finite-dimensional, and it leaves the interpretive core of quantum mechanics untouched: geometry alone does not select a pointer basis, does not guarantee that a given Hamiltonian drives the system into typical regions of the accessible subspace, and does not turn an improper mixture into a proper one. It merely supplies the vast Hilbert-space capacity that makes decoherence so overwhelmingly effective for all practical purposes.
\end{abstract}

\maketitle

\section{Introduction}

The quantum measurement problem asks how the continuous and reversible Schr\"odinger evolution gives rise to the discrete, irreversible classical world we observe~\cite{schlosshauer-2007,schlosshauer-2019}. Operationally, the question reduces to why the off-diagonal interference terms of a system's reduced density matrix are unobservable on macroscopic scales.

The urgency of this question was forcefully articulated by Schr\"odinger in a 1952 colloquium~\cite{schroedinger-interpretation}. He noted that if the wave equation is universally valid, then alternative macroscopic outcomes must all ``really happen simultaneously.'' Schr\"odinger argued that if nature actually took this form, we should find our surroundings rapidly blurring together into a ``quagmire, or sort of a featureless jelly''---a variant of his cat metaphor. For example, a watch left unobserved in a drawer overnight should dissolve into a smeared probability distribution. This ``jellification'' argument poses a sharp, physical question: if unitary quantum mechanics holds at all scales, what fundamental mechanism prevents the simultaneous alternatives of the macroscopic world from blurring into an unobservable, featureless smear?

Bell coined ``FAPP'' (\emph{for all practical purposes}) to express skepticism toward explanations that rest on the practical impossibility of observing macroscopic interference rather than on a fundamental physical mechanism~\cite{bell-a}. Modern decoherence theory~\cite{joos1985emergence,RevModPhys.75.715,joos2003decoherence,schlosshauer-2007} answers this by showing that interaction with a macroscopic environment dynamically drives the off-diagonal terms toward zero in a preferred (pointer) basis selected by the system--environment coupling.

We shall examine a complementary, purely \emph{geometric}
contribution to the same suppression.  The relevant tools are the
exponential packing of approximately orthogonal unit vectors and the
concentration of measure on high-dimensional spheres
\cite{kainen1993quasiorthogonal,tao2012topics,vershynin2018high,ledoux2001concentration,milman1986asymptotic}.  These results are
standard in high-dimensional geometry, but their direct application to
the environmental records appearing in decoherence is, to our knowledge,
underemphasized.

Our central claim is modest. High-dimensional Hilbert-space geometry guarantees an enormous reservoir of nearly non-interfering pure states. This is a kinematic fact independent of any Hamiltonian. Whether actual environmental branches occupy this reservoir is a separate dynamical question. In ordinary decoherence theory, the system--environment interaction selects a pointer basis and correlates different pointer values with different environmental records. The geometric estimates below apply once those records are sufficiently typical in the accessible environmental subspace.

We therefore separate two ingredients that are often bundled together. The dynamical ingredient is the production of environmental records by the Hamiltonian. The kinematic ingredient is the concentration of measure on high-dimensional spheres, which makes generic records almost mutually orthogonal. The latter sharply quantifies the FAPP suppression of macroscopic interference, while leaving untouched the deeper interpretive question of whether an improper mixture has become a proper one.


\section{Random Quantum States, Typical Overlaps, and Exponential Packing}
\label{sec:random-overlaps-packing}

\subsection*{Heuristic argument}

The main geometric point can be understood before any probability formalism is
introduced.  A normalized quantum state
\begin{equation}
    \ket{\psi}
    =
    (\psi_1,\ldots,\psi_d)
    \in \mathbb C^d
\end{equation}
satisfies
\begin{equation}
    |\psi_1|^2+\cdots+|\psi_d|^2=1 .
\end{equation}
If the state is chosen without any preferred direction, this unit probability
weight is spread over $d$ complex coordinates.  Thus a typical coordinate
weight should be of order $1/d$.  Now fix a pure state $\ket{\phi}$.  By
unitary invariance one may choose coordinates so that $\ket{\phi}=\ket{e_1}$.
The Born probability for a system prepared in $\ket{\psi}$ to pass the yes--no
test associated with $\ket{\phi}$ is then
\begin{equation}
    |\langle\phi|\psi\rangle|^2=|\psi_1|^2 .
\end{equation}
Hence the transition probability from a random pure state to a fixed pure state
is expected to be of order $1/d$.  A transition probability larger than some
fixed threshold $\eps>0$ would require one coordinate to carry an anomalously
large fraction of the total norm.  In high dimension this is exponentially
unlikely.  The exact calculation below makes this heuristic precise:
\begin{equation}
    \PP\{ |\langle\phi|\psi\rangle|^2\geq\eps\}
    =
    (1-\eps)^{d-1}
    \leq
    \exp[-(d-1)\eps].
    \label{eq:heuristic-tail}
\end{equation}
This exact overlap tail is the source of both typical quasi-orthogonality and
the exponential packing estimate used below.

\subsection*{Notation and elementary distributions}

The symbol $\PP$ denotes the probability of an event, and $\E$ denotes
expectation with respect to the random choice of state vectors.  A Haar-random
unit vector in $\mathbb C^d$ means a unit vector distributed according to the
unique unitarily invariant probability measure on the unit sphere: applying
any unitary $U\in U(d)$ leaves its distribution unchanged.  In the real
comparison below, Haar-random means uniformly distributed with respect to the
orthogonally invariant probability measure on the real sphere
$S^{d-1}\subset\mathbb R^d$.

The notation $N(0,1)$ denotes the real normal distribution with mean zero and
variance one, with density
\begin{equation}
    (2\pi)^{-1/2}\exp(-x^2/2),
    \qquad x\in\mathbb R .
\end{equation}
Thus $X,Y\sim N(0,1)$ means that $X$ and $Y$ are independent standard real
Gaussian random variables.  A standard complex Gaussian variable is
\begin{equation}
    Z=\frac{X+iY}{\sqrt2}.
    \label{eq:standard-complex-gaussian}
\end{equation}
The factor $1/\sqrt2$ normalizes the second moment:
\begin{equation}
    \E |Z|^2
    =
    \frac{\E X^2+\E Y^2}{2}
    =1 .
\end{equation}
Equivalently, $Z$ has density $\pi^{-1}e^{-|z|^2}$ on the complex plane.  A
standard complex Gaussian vector $Z=(Z_1,\ldots,Z_d)\in\mathbb C^d$ has
independent standard complex Gaussian components.

We shall use a few elementary distributional facts.  A chi-square random
variable with $k$ degrees of freedom, denoted $\chi^2_k$, is the sum of the
squares of $k$ independent $N(0,1)$ variables.  Hence, if
$Z=(X+iY)/\sqrt2$ is a standard complex Gaussian, then
\begin{equation}
    |Z|^2
    =
    \frac{X^2+Y^2}{2},
    \qquad
    X^2+Y^2\sim\chi^2_2 .
\end{equation}
Geometrically, $X^2+Y^2$ is the squared radius of a two-dimensional isotropic
Gaussian point.  Its density is
\begin{equation}
    \frac12 e^{-s/2},
    \qquad s\geq0,
\end{equation}
so $(X^2+Y^2)/2$ has density $e^{-t}$ on $t\geq0$, i.e. the exponential
distribution of mean one.  We write this as
\begin{equation}
    |Z|^2\sim\operatorname{Exp}(1)
    =
    \operatorname{Gamma}(1,1).
\end{equation}
For different coordinates $j$, the variables $|Z_j|^2$ are independent because
the Gaussian components $Z_j$ are independent and $|Z_j|^2$ depends only on the
$j$th component.

We use the shape--scale convention for gamma distributions:
$\operatorname{Gamma}(\alpha,\theta)$ has density proportional to
$x^{\alpha-1}e^{-x/\theta}$ on $x>0$, with shape $\alpha$ and scale
$\theta$.  Thus $\operatorname{Gamma}(\alpha,1)$ has unit scale.  The beta
distribution $\operatorname{Beta}(a,b)$ has density
\begin{equation}
    \frac{1}{B(a,b)}t^{a-1}(1-t)^{b-1},
    \qquad 0<t<1,
\end{equation}
where $B(a,b)$ is the beta function, i.e. the normalizing integral.  We use the
standard beta-ratio fact: if $U\sim\operatorname{Gamma}(a,\theta)$ and
$V\sim\operatorname{Gamma}(b,\theta)$ are independent with the same scale
$\theta$, then
\begin{equation}
    \frac{U}{U+V}\sim\operatorname{Beta}(a,b).
\end{equation}
The scale parameter cancels in the ratio.

\subsection*{Exact complex overlap distribution}

The point is directly quantum-mechanical.  If a system is prepared in
$\ket{\psi}$ and tested against the pure state $\ket{\phi}$, the Born
probability of the corresponding yes-outcome is
\begin{equation}
    \bra{\psi}\bigl(\ket{\phi}\bra{\phi}\bigr)\ket{\psi}
    =
    |\langle\phi|\psi\rangle|^2 .
    \label{eq:born-transition-probability}
\end{equation}
Thus quasi-orthogonality of random pure states is equivalently a statement
about the typical smallness of transition probabilities.

\begin{lemma}[Complex random transition probabilities]
\label{lem:complex-transition-probability}
Let $\ket{\psi},\ket{\phi}\in\mathbb C^d$ be independently chosen
Haar-random unit vectors, with $d\geq2$. Then, for $0\leq \eps\leq1$,
\begin{equation}
    \PP\{ |\langle\phi|\psi\rangle|^2\geq \eps\}
    =
    (1-\eps)^{d-1} .
    \label{eq:overlapExact}
\end{equation}
Equivalently, for $0\leq\delta\leq1$,
\begin{equation}
    \PP\{ |\langle\phi|\psi\rangle|\geq \delta\}
    =
    (1-\delta^2)^{d-1}
    \leq
    \exp[-(d-1)\delta^2] .
    \label{eq:complex-overlap-tail-delta}
\end{equation}
\end{lemma}

\begin{proof}
By unitary invariance of Haar measure, we may fix
$\ket{\phi}=\ket{e_1}$, the first vector of an orthonormal basis.  Indeed,
conditioning on any value of $\ket{\phi}$, a unitary sending that vector to
$\ket{e_1}$ leaves the distribution of the independently chosen $\ket{\psi}$
Haar-random.  Thus averaging over two independent random vectors has the same
overlap distribution as fixing one vector and averaging over the other.  It
remains to determine the distribution of $|\psi_1|^2$, where
$\ket{\psi}=(\psi_1,\ldots,\psi_d)$ is Haar-random on the unit sphere of
$\mathbb C^d$.

A Haar-random unit vector can be generated by normalizing a complex Gaussian
vector.  Let
\begin{equation}
    Z=(Z_1,\ldots,Z_d),
\end{equation}
where the $Z_j$ are independent standard complex Gaussian variables, and set
\begin{equation}
    \ket{\psi}=\frac{Z}{\|Z\|} .
\end{equation}
This construction gives the Haar distribution because the joint density of
$Z$ is proportional to
\begin{equation}
    \exp(-\|Z\|^2),
\end{equation}
and therefore depends only on the norm $\|Z\|$.  It is invariant under every
unitary transformation $U\in U(d)$: the random vectors $Z$ and $UZ$ have the
same distribution. Consequently $Z/\|Z\|$ and $UZ/\|Z\|$ have the same
distribution on the unit sphere.  By uniqueness of the unitarily invariant
probability measure on that sphere, this distribution is precisely Haar
measure.  Then
\begin{equation}
    |\psi_1|^2
    =
    \frac{|Z_1|^2}{|Z_1|^2+\cdots+|Z_d|^2} .
\end{equation}
By Eq.~\eqref{eq:standard-complex-gaussian}, each component can be written as
$Z_j=(X_j+iY_j)/\sqrt2$, with $X_j,Y_j$ independent $N(0,1)$ variables.
Therefore
\begin{equation}
    |Z_j|^2
    =
    \frac{X_j^2+Y_j^2}{2}.
\end{equation}
Since $X_j^2+Y_j^2\sim\chi^2_2$, and $\chi^2_2/2$ has the exponential
distribution of mean one, $|Z_j|^2\sim\operatorname{Exp}(1)$, equivalently
$|Z_j|^2\sim\operatorname{Gamma}(1,1)$.  The independence of the $|Z_j|^2$
follows from the independence of the pairs $(X_j,Y_j)$.  To avoid overloading
notation, write
\begin{equation}
    E_j=|Z_j|^2,
\end{equation}
with the $E_j$ independent exponential random variables of mean one.  Let
\begin{equation}
    W=E_2+\cdots+E_d .
\end{equation}
Then $E_1$ and $W$ are independent, with
$E_1\sim\operatorname{Gamma}(1,1)$ and
$W\sim\operatorname{Gamma}(d-1,1)$. Therefore
\begin{equation}
    |\psi_1|^2
    =
    \frac{E_1}{E_1+W}
    \sim
    \operatorname{Beta}(1,d-1).
    \label{eq:complex-beta-distribution}
\end{equation}
This gives a useful symmetry with the real case below: the complex transition
probability is beta distributed with parameters $(1,d-1)$, whereas the real
case has parameters $(1/2,(d-1)/2)$.

For completeness, and also to make the tail explicit, we compute the same
probability directly. For $0\leq\eps<1$,
\begin{align}
 \PP\{|\psi_1|^2\geq\eps\}
 &=
 \PP\left\{\frac{E_1}{E_1+W}\geq\eps\right\}
 \nonumber\\
 &=
 \PP\left\{E_1\geq \frac{\eps}{1-\eps}W\right\}
 \nonumber\\
 &=
 \E\left[
        \exp\left(-\frac{\eps}{1-\eps}W\right)
      \right].
\end{align}
The last equality is just conditioning on $W$: since $E_1$ is independent of
$W$ and has exponential tail $\PP\{E_1\geq t\}=e^{-t}$, one has
\begin{equation}
    \PP\{E_1\geq sW\mid W=w\}=e^{-sw},
\end{equation}
and averaging over $W$ gives $\E(e^{-sW})$.  Since $W$ is the sum of $d-1$
independent exponential variables, $W\sim\operatorname{Gamma}(d-1,1)$, and
its Laplace transform is
\begin{equation}
    \E(e^{-sW})=(1+s)^{-(d-1)} .
\end{equation}
With $s=\eps/(1-\eps)$, this gives
\begin{equation}
    \PP\{|\psi_1|^2\geq\eps\}
    =
    (1-\eps)^{d-1} .
\end{equation}
At the endpoint $\eps=1$, the same formula gives zero; equivalently, a
Haar-random vector is parallel to the fixed vector with probability zero for
$d\geq2$.  The same formula also follows immediately from the cumulative
distribution function of $\operatorname{Beta}(1,d-1)$, since
\begin{equation}
    \PP\{T\geq\eps\}=(1-\eps)^{d-1},
    \qquad T\sim\operatorname{Beta}(1,d-1).
\end{equation}
Since $|\langle\phi|\psi\rangle|^2=|\psi_1|^2$,
Eq.~\eqref{eq:overlapExact} follows.  Eq.~\eqref{eq:complex-overlap-tail-delta}
follows by putting $\eps=\delta^2$, because $\delta$ is the amplitude-overlap
threshold whereas $\eps$ is the corresponding Born-probability threshold, and
then using $1-x\leq e^{-x}$.
\end{proof}

Several consequences are immediate.  The mean transition probability is
\begin{equation}
    \E |\langle\phi|\psi\rangle|^2
    =
    \frac{1}{d},
    \label{eq:mean-transition-probability}
\end{equation}
which also follows from symmetry, since $\sum_{j=1}^d|\psi_j|^2=1$. Moreover,
for every fixed threshold $\eps>0$,
\begin{equation}
    \PP\{ |\langle\phi|\psi\rangle|^2\geq\eps\}
    \leq
    \exp[-(d-1)\eps] .
    \label{eq:large-transition-probability-suppressed}
\end{equation}
Thus large transition probabilities between independently chosen random pure
states are exponentially suppressed in the Hilbert-space dimension. In this
precise sense, two random pure states are quasi-orthogonal with probability
tending to one exponentially fast with $d$.

It is useful to separate two scales.  The typical Born probability
$|\langle\phi|\psi\rangle|^2$ is of order $d^{-1}$, while the typical amplitude
$|\langle\phi|\psi\rangle|$ is of order $d^{-1/2}$. Thus random pure states are
not exactly orthogonal, but their mutual transition probabilities are normally
microscopic in large Hilbert spaces.

\subsection*{Exponential packing consequence}

In a $d$-dimensional Hilbert space $\HH_d$, the maximum number of exactly
mutually orthogonal unit vectors is $d$.  We relax exact orthogonality as
follows. Two unit vectors $\ket{\psi_i},\ket{\psi_j}\in\HH_d$ are called
$\eps$-quasi-orthogonal if
\begin{equation}
    |\langle\psi_i|\psi_j\rangle|^2 \le \eps .
\end{equation}
A family $\{\ket{\psi_i}\}_{i=1}^{M}$ is $\eps$-quasi-orthogonal if this
inequality holds for every pair $i\ne j$.

Let $M_\eps(d)$ denote the largest possible cardinality of such a family.  The
exact overlap tail above gives the following random-coding lower bound:
\begin{equation}
    M_\eps(d)
    \;\ge\;
    \left\lfloor
    \exp\!\left[
        \frac{d-1}{2}\bigl(-\log(1-\eps)\bigr)
    \right]
    \right\rfloor .
    \label{eq:packing-lower}
\end{equation}
Indeed, sample $M$ unit vectors $\ket{\psi_1},\dots,\ket{\psi_M}\in\HH_d$
independently from Haar measure.  For any pair $i\ne j$, Eq.~\eqref{eq:overlapExact}
gives
\begin{equation}
\PP\bigl(|\braket{\psi_i}{\psi_j}|^2 \ge \eps\bigr)
= (1-\eps)^{d-1}.
\end{equation}
By a union bound over the $\binom{M}{2}<M^2/2$ pairs,
\begin{equation}
\PP\bigl(\exists\, i\ne j:
 |\braket{\psi_i}{\psi_j}|^2 \ge \eps\bigr)
<
\frac12 M^2(1-\eps)^{d-1}.
\end{equation}
Choosing
\begin{equation}
M=
\left\lfloor
\exp\!\left[
        \frac{d-1}{2}\bigl(-\log(1-\eps)\bigr)
    \right]
\right\rfloor
\end{equation}
gives
\begin{equation}
\frac12 M^2(1-\eps)^{d-1}\leq \frac12<1,
\end{equation}
so the bad event cannot have probability one. Hence there exists at least
one $\eps$-quasi-orthogonal family of cardinality $M$, proving
Eq.~\eqref{eq:packing-lower}.

For small fixed $\eps$, since $-\log(1-\eps)=\eps+O(\eps^2)$, this implies
\begin{equation}
    M_\eps(d)\ge \exp(c\,\eps d)
\end{equation}
for some constant $c>0$, fixed $\eps>0$, and sufficiently large $d$.  The
statement should be understood in the fixed-$\eps$, large-$d$ regime; no bound
of the form $M_\eps(d)\le \exp(C\eps d)$ can hold uniformly down to
$\eps=O(1/d)$, because an orthonormal basis already gives
$M_\eps(d)\ge d$ for every $\eps\ge0$.  Standard spherical-code bounds also
give exponential-in-$d$ upper bounds for fixed $\eps>0$~\cite{kabatiansky1978bounds,conway1999sphere}, but the precise optimal exponent is not needed here. We therefore use only the random-coding lower bound.

For a system of $n$ qubits, $d=2^n$. Hence, for fixed $\eps>0$,
\begin{equation}
    M_\eps(2^n)
    \;\ge\;
    \exp\!\bigl(c\,\eps\,2^n\bigr),
\end{equation}
which is doubly exponential in the number $n$ of qubits. Thus the Hilbert
space of a macroscopic environment has vastly more than enough kinematic room
to accommodate an enormous number of mutually quasi-orthogonal environmental
records.

It is important not to overinterpret this packing statement.  The bound concerns
pairwise transition probabilities between pure states. It does not imply that a
$d$-dimensional environment can store exponentially many reliably distinguishable
classical records. If $M\gg d$, the packed states are necessarily linearly
dependent, and pairwise small overlaps alone do not guarantee a measurement that
recovers the corresponding labels with high average success probability. The
relevance for decoherence is more modest and more appropriate: for a fixed,
moderate, or dynamically selected set of macroscopic alternatives, the
suppression of pairwise overlaps is the quantity that controls the off-diagonal
terms of the reduced state. In that setting, high-dimensional Hilbert space
provides ample kinematic room for environmental branch states to become
effectively orthogonal.

This statement is closely related in spirit to the concentration phenomenon
underlying the Johnson--Lindenstrauss lemma, but the emphasis is different.
Johnson--Lindenstrauss is a finite-set embedding result: $N$ vectors can be
randomly projected to $O(\eps^{-2}\log N)$ dimensions while preserving pairwise
distances up to relative error $\eps$.  The present argument uses the
complementary viewpoint.  In a fixed high-dimensional Hilbert space,
independently chosen pure states already have Born transition probabilities of
order $1/d$, and large overlaps are exponentially suppressed.  Thus
Johnson--Lindenstrauss concerns the preservation of geometry under random
dimensional reduction, whereas the present use concerns the generic
near-orthogonality already available in high dimension.

\subsection*{Real comparison}
\label{subsec:real-case}

The complex case is the natural one for quantum mechanics.  For comparison,
the real case has the same concentration behaviour, but a slightly less
elementary closed form.  Let $u,v\in\mathbb R^d$ be independent Haar-random
unit vectors on the real sphere $S^{d-1}$.  By orthogonal invariance, fix
$v=e_1$.  A Haar-random real unit vector is generated analogously by
normalizing a standard real Gaussian vector: if $G=(G_1,\ldots,G_d)$, where
the $G_j$ are independent $N(0,1)$ variables, then $G/\|G\|$ is invariant
under all orthogonal transformations and hence is uniform on $S^{d-1}$. Thus,
for $u=G/\|G\|$,
\begin{equation}
    |\langle v,u\rangle|^2
    =
    \frac{G_1^2}{G_1^2+\cdots+G_d^2} .
\end{equation}
Here $G_1^2\sim\operatorname{Gamma}(1/2,2)$, and
$G_2^2+\cdots+G_d^2\sim\operatorname{Gamma}((d-1)/2,2)$, in the same
shape--scale convention.  The scale is now $2$ rather than $1$, but it is
common to numerator and denominator and therefore cancels in the ratio.  The
beta-ratio rule gives
\begin{equation}
    |\langle v,u\rangle|^2
    \sim
    \operatorname{Beta}\left(\frac12,\frac{d-1}{2}\right).
    \label{eq:real-beta-distribution}
\end{equation}
Equivalently, its density on $(0,1)$ is
\begin{equation}
    f_d(t)
    =
    \frac{t^{-1/2}(1-t)^{(d-3)/2}}
         {B\left(\frac12,\frac{d-1}{2}\right)} .
    \label{eq:real-beta-density}
\end{equation}
Thus
\begin{equation}
    \PP\{|\langle v,u\rangle|^2\geq\eps\}
    =
    \int_\eps^1
    \frac{t^{-1/2}(1-t)^{(d-3)/2}}
         {B\left(\frac12,\frac{d-1}{2}\right)}\,dt .
    \label{eq:real-tail-integral}
\end{equation}
For fixed $\eps>0$, this again decays exponentially in $d$, because the factor
$(1-t)^{(d-3)/2}$ is exponentially small throughout the integration range
$t\geq\eps$. More explicitly,
\begin{align}
 \int_\eps^1 t^{-1/2}(1-t)^{(d-3)/2}\,dt
 &\leq
 \eps^{-1/2}
 \int_\eps^1 (1-t)^{(d-3)/2}\,dt
 \nonumber\\
 &=
 \frac{2\eps^{-1/2}}{d-1}
 (1-\eps)^{(d-1)/2} .
\end{align}
Furthermore, $B(1/2,(d-1)/2)^{-1}=O(\sqrt d)$ by Stirling's formula. Hence,
for each fixed $\eps\in(0,1)$,
\begin{align}
    \PP\{|\langle v,u\rangle|^2\geq\eps\}
    &\leq
    C_\eps d^{-1/2}(1-\eps)^{(d-1)/2}
    \nonumber\\
    &\leq
    C_\eps d^{-1/2}
    \exp\left[-\frac{(d-1)\eps}{2}\right]
\end{align}
for a constant $C_\eps$ independent of $d$.  The exponent is roughly half the
complex exponent, reflecting the fact that $\mathbb C^d$ has twice as many
real degrees of freedom as $\mathbb R^d$.  The qualitative high-dimensional
conclusion is nevertheless the same: a random real unit vector is almost
orthogonal to an independently chosen one with probability tending to one
exponentially fast in the dimension.

\subsection*{Relation to L\'evy's lemma}
\label{subsec:relation-to-levy}

The preceding argument is self-contained and uses no spherical isoperimetric
black box.  It exploits the special observable relevant to pure-state quantum
transition probabilities, namely
\begin{equation}
    \ket{\psi}\longmapsto |\langle\phi|\psi\rangle|^2 .
\end{equation}
L\'evy's lemma is the corresponding general concentration theorem: it replaces
this particular transition probability by an arbitrary Lipschitz observable on
the high-dimensional unit sphere.  In that broader setting one obtains bounds
of the schematic form
\begin{equation}
    \PP\{|f-\text{typical value}|\geq\eta\}
    \lesssim
    \exp\left[-c\,\frac{d\eta^2}{L^2}\right],
\end{equation}
where $L$ is the Lipschitz constant of $f$.  Proving this general statement
requires the usual concentration-of-measure machinery, for example spherical
isoperimetry.

The transition-probability observable itself is Lipschitz. If
$P_\phi=\ket{\phi}\bra{\phi}$ and
$f(\psi)=\bra{\psi}P_\phi\ket{\psi}=|\langle\phi|\psi\rangle|^2$, then for
two unit vectors $\psi,\psi'$,
\begin{equation}
    |f(\psi)-f(\psi')|
    \leq
    2\,\|\psi-\psi'\| .
\end{equation}
Thus $f$ has Lipschitz constant at most $2$ with respect to Euclidean chordal
distance, and hence also with respect to geodesic distance up to the usual
constant convention. Applying L\'evy's lemma to this particular observable
would therefore prove concentration, but with a generic Gaussian tail in the
deviation parameter.  The exact beta-distribution calculation above is sharper
for transition probabilities: for fixed $\eps$, it gives an exponent linear in
$\eps$, namely $\exp[-(d-1)\eps]$, rather than the weaker generic
Lipschitz-type behaviour quadratic in the deviation. For the quantum
quasi-orthogonality claim above, no black-box input is needed: the exact
distribution of the Born transition probability already gives the desired
exponential suppression.

\section{Application to the Measurement Process}

Consider a system $\SSys$ coupled to an apparatus--environment $\EE$ initialised in $\ket{E_0}$:
\begin{equation}
    \ket{\Psi(0)} = \Bigl(\sum_i c_i \ket{s_i}\Bigr)\otimes\ket{E_0}.
\end{equation}
A measurement-like unitary correlates pointer outcomes with environmental records:
\begin{equation}
    \ket{\Psi(t)} = \sum_i c_i \ket{s_i}\otimes\ket{E_i(t)}.
\end{equation}
The reduced density matrix on $\SSys$ is
\begin{equation}
    \rho_{\SSys}(t) = \sum_{i,j} c_i c_j^{*}\,\braket{E_j(t)}{E_i(t)}\,\ket{s_i} \bra{s_j}.
\end{equation}
Off-diagonal coherences are weighted by the environmental overlaps
$\braket{E_j(t)}{E_i(t)}$, and are suppressed in modulus by at most
$\sqrt{\varepsilon}$ whenever the branch states are $\varepsilon$-quasi-orthogonal.

\subsection*{The kinematic claim.}
Throughout this section we adopt natural units $k_B=\hbar=1$, so that entropies are dimensionless and inverse temperatures have units of time. \emph{If} the branch states $\{\ket{E_i(t)}\}_{i\ne j}$ are well-modeled as a typical pair on the relevant high-dimensional subspace of $\HH_{\EE}$, then
\begin{equation}
    \E\bigl[\,|\braket{E_j(t)}{E_i(t)}|^2\,\bigr] \;=\; \frac{1}{d_{\rm eff}},
\end{equation}
with $d_{\rm eff}$ the dimension of the subspace explored by the dynamics. For a generic chaotic many-body environment the accessible subspace is the microcanonical shell, so
\begin{equation}
    d_{\rm eff} \;\sim\; e^{S},
    \qquad S \sim n,
    \label{eq:dEntropy}
\end{equation}
with $S$ the (dimensionless) microcanonical entropy of the relevant environmental sector and $n$ the number of constituent degrees of freedom; in conventional units the right-hand sides of~\eqref{eq:dEntropy} read $e^{S/k_B}$ and $nk_B$, respectively. Equation~\eqref{eq:dEntropy} ties the geometric parameter $d_{\rm eff}$ directly to standard thermodynamics: the mean-square overlap is exponential in the entropy of the environment, so the typical overlap amplitude is of order $e^{-S/2}$.

The relevant dimension is not necessarily the full Hilbert-space dimension of
$\HH_{\EE}$, but the dimension of the subspace $\HH_{\rm acc}$ accessible to
the dynamics from the initial state on the timescale considered. Energy
conservation, conserved charges, and finite-time evolution can all restrict the
accessible subspace. For a many-body environment with Hamiltonian $H_{\EE}$ and
an initial state localized within an energy window $[E,E+\Delta E]$, the
accessible subspace is often well approximated by the microcanonical shell,
with
\begin{equation}
    d_{\rm eff}=\dim\HH_{\rm acc}\sim e^{S(E)/k_B}.
    \label{eq:microcanonical-deff}
\end{equation}
In the natural units used here, this is simply $d_{\rm eff}\sim e^{S(E)}$.
For chaotic many-body dynamics with bounded local interactions, $S(E)$ is
extensive in the number of constituents, so $d_{\rm eff}$ remains exponential
in system size. For integrable or many-body-localized environments, additional
conserved quantities can drastically reduce $d_{\rm eff}$; the geometric
reservoir is then correspondingly smaller.

This is vastly stronger than needed to render coherences operationally invisible on any timescale shorter than the Poincar\'e recurrence time, which for a generic chaotic many-body system is astronomically large---heuristic arguments suggest a scaling at least as $\exp(d_{\rm eff})\sim\exp(e^{S})$, but the precise value is model-dependent and irrelevant to our argument~\cite{peres1982recurrence,wallace2015recurrence}.


\subsection*{Minimal dynamical input.}

The argument is kinematic rather than dynamical. It assumes that the interaction
Hamiltonian has already selected the relevant pointer alternatives and that the
associated environmental branch states explore a sufficiently large accessible
subspace in a typical manner. It does not by itself derive the pointer basis,
nor does it prove thermalization or equilibration. Rather, it explains why, once
distinct branches are driven into a large environmental Hilbert space, their
mutual overlaps are generically expected to be extremely small.

In a measurement-like interaction, after a pointer basis has been selected, one may write schematically
\begin{equation}
    \ket{E_i(t)}
    =
    U_i(t)\ket{E_0},
\end{equation}
where the conditional unitary $U_i(t)$ depends on the pointer value $i$. The geometric argument does not prove that the vectors $U_i(t)\ket{E_0}$ are Haar-random. The required dynamical assumption is only the weaker pair-typicality condition that, for $i\ne j$, their overlap statistics are close to those of two typical vectors in the accessible subspace $\mathcal H_{\rm acc}$:
\begin{equation}
    |\langle E_j(t)|E_i(t)\rangle|^2
    \sim
    d_{\rm eff}^{-1},
    \qquad
    d_{\rm eff}=\dim\mathcal H_{\rm acc}.
\end{equation}

Chaotic, nonintegrable environments provide the usual physical reason to expect such typicality after a suitable relaxation or decoherence time. Integrable, many-body-localized, weakly coupled, or strongly constrained environments may fail to satisfy it. Thus the geometry supplies the available reservoir of nearly orthogonal records, while the Hamiltonian must still populate that reservoir. This assumption is in the same broad family as the typicality assumptions used in canonical typicality and quantum thermalization~\cite{popescu2006entanglement,goldstein2006canonical,reimann2008foundation,rigol2008thermalization,dalessio2016quantum}.

\subsection*{Limitations.}

Two limitations should be kept explicit. First, quasi-orthogonality
does not select the pointer basis. The preferred basis is determined
by the structure of the system--environment interaction, not by
high-dimensional geometry alone.

 Second, small off-diagonal terms do
not by themselves turn an improper mixture into a proper one. The
geometric estimates therefore strengthen the FAPP suppression of
interference, but they do not resolve the interpretive core of the
measurement problem.

Third, pairwise suppression should not be conflated with arbitrary classical
readout capacity. If the number of alternatives is so large that the full Gram
matrix is densely populated by small but numerous off-diagonal terms, pairwise
bounds alone need not control all global decoding or distinguishability
questions. The decoherence application therefore concerns fixed, moderate, or
dynamically selected families of macroscopic alternatives, not an exponentially
large alphabet of perfectly readable records.

A strict, rather than merely FAPP, breakdown
of unitary equivalence between macroscopic outcomes can be obtained
in the infinite tensor product limit through sectorization and
factorization~\cite{vonNeumann1939,hepp-1972,van-den-bossche-2023-a,svozil-2025-u},
of which the present finite-dimensional, geometric estimates may be
viewed as a quantitative companion: exponentially small (rather than
exactly vanishing) overlaps between macroscopically distinct
environmental records.


\subsection*{Schr\"odinger's ``jellification'' worry.}
\label{sec:jelly}

Schr\"odinger's worry was that, if unitary quantum mechanics applies universally, macroscopically distinct alternatives should blur into a single ``featureless jelly.'' In the present language, such blurring would require appreciable overlaps between the environmental records correlated with different macroscopic alternatives. Under the pair-typicality assumption above, however,
\[
    |\langle E_j(t)|E_i(t)\rangle|^2
    \sim
    d_{\rm eff}^{-1}
    \sim
    e^{-S},
\]
so the typical overlap amplitude is of order $e^{-S/2}$. For a macroscopic environment this is FAPP indistinguishable from zero. The branches do not blur because the accessible Hilbert space contains an enormous supply of nearly orthogonal directions. This answers the jellification worry operationally, while leaving the deeper interpretive questions untouched.

\section{Discussion}

The geometric perspective separates two ingredients in decoherence. The first is dynamical: the interaction Hamiltonian selects a pointer basis and correlates distinct pointer values with distinct environmental records. The second is kinematic: once those records behave like typical vectors in a high-dimensional accessible subspace, their mutual overlaps are automatically tiny.

For an accessible subspace of dimension $d_{\rm eff}\sim e^S$, the typical squared overlap is $d_{\rm eff}^{-1}\sim e^{-S}$, and the typical overlap amplitude is $e^{-S/2}$. This suppression is not a new dynamical mechanism; it is the high-dimensional geometry underlying the effectiveness of environmental decoherence. The geometry explains why there is ample Hilbert-space room for many macroscopic records to coexist without appreciable interference, but it does not explain why a particular Hamiltonian produces those records, nor does it select a preferred basis.

This distinction also clarifies where the framework can fail. If the environment is integrable, many-body localized, weakly coupled, or restricted by additional conserved quantities, the branch states may not be typical in the relevant subspace. In that case the geometric reservoir exists but is not dynamically populated.

The argument is close in spirit to canonical typicality, where concentration of measure is used to show that most pure states of a large composite system have nearly thermal reduced states on small subsystems~\cite{popescu2006entanglement,goldstein2006canonical,reimann2008foundation}. Here the object of interest is different: the pairwise overlap between distinct environmental records. Infinite-tensor-product approaches to measurement obtain exact orthogonality between macroscopic sectors in the limit $n=\infty$; the present analysis can be read as the finite-$n$ version, where exact sectorization is replaced by exponentially small overlap.

\section{Conclusion}

We have isolated a kinematic contribution to the suppression of macroscopic interference. In a $d$-dimensional Hilbert space, random pure states are almost orthogonal: for fixed $\ket{\phi}$ and Haar-random $\ket{\psi}$,
\[
    \mathbb E|\langle\phi|\psi\rangle|^2=\frac1d,
    \qquad
    \mathbb P\!\left(|\langle\phi|\psi\rangle|^2\ge\eps\right)
    =(1-\eps)^{d-1}.
\]
Moreover, for fixed $\eps>0$, random coding produces $\eps$-quasi-orthogonal families of size at least $\exp(c\eps d)$. For an environmental subspace with $d_{\rm eff}\sim e^S$, this means that typical squared branch overlaps are of order $e^{-S}$.

These facts do not replace decoherence dynamics. The Hamiltonian must still select the pointer basis and drive environmental records into typical regions of the accessible subspace. Nor do the estimates convert an improper mixture into a proper one. What the geometry does provide is a quantitative explanation of why, once environmental records have formed, macroscopic alternatives can coexist without producing Schr\"odinger's ``featureless jelly'': high-dimensional Hilbert space supplies an enormous reservoir of nearly non-interfering directions.

\begin{acknowledgments}
This text was partially created and revised with assistance from large language models.
This research was funded in whole or in part by the Austrian Science Fund (FWF) Grant DOI: 10.55776/PIN5424624.
The author acknowledges TU Wien Bibliothek for financial support through its Open Access Funding Programme.
\end{acknowledgments}

\bibliography{svozil}

\end{document}
